\section*{Roadmap}
\addcontentsline{toc}{section}{Roadmap}

\noindent\motiv{\vet{\extraonly{\revd{The lattice appears in Chapter~\ref{sec:nets} as an example;
its propagators (Chapter~\ref{sec:lattice}) come before the general net
(Chapter~\ref{sec:phasespace}), and the lattice net is completed in
\reve{Chapter~\ref{sec:latticenet}}\revg{, apart from primitivity (Chapter~\ref{sec:simple})}.}}\dimockonly{\revm{This
edition contains only what the verification of the axioms of~\cite{dimock} for the free scalar field
needs; the full notes add the lattice field, the net defined through Cauchy data, Huygens' principle
and some illustrations.}}
\revm{Chapter~\ref{sec:nets} states the axioms for any set of regions carrying inclusion,
orthogonality and dependence, and reduces them to properties of generators
(Proposition~\ref{prop:generated}); Chapter~\ref{sec:weyl} constructs the Weyl operators and the CCR
algebra.} Figure~\ref{fig:roadmap} shows which chapters use which.}}

\ifdimockedition
\begin{figure}[ht]
\centering
\begin{tikzpicture}[
  font=\small,
  done/.style={draw, rounded corners=2pt, fill=Blue!6, text width=3.5cm, align=center,
    minimum height=1.05cm, inner sep=3pt},
  dep/.style={-{Stealth[length=2mm]}, thick},
  node distance=0.9cm and 0.7cm]
  \node[done] (c1) {\ref*{sec:nets}\quad Causal indices and nets of local algebras};
  \node[done, right=of c1] (c2) {\ref*{sec:weyl}\quad The Weyl algebra};
  \node[done, below=of c1] (c4) {\ref*{sec:phasespace}\quad The net of a classical phase space};
  \node[done, below=of c2] (c6) {\ref*{sec:simple}\quad Simplicity and primitivity of the Weyl algebra};
  \node[done, below=of c4] (c7) {\ref*{sec:green}\quad The net of a Green-hyperbolic operator};
  \node[done, below=of c6] (c8) {\ref*{sec:minkowski}\quad The free massless\\ field on Minkowski space};
  \node[done, below=of c8] (c9) {\ref*{sec:mass}\quad The Klein--Gordon\\ field on Minkowski space};
  \draw[dep] (c1) -- (c2);
  \draw[dep] (c2) -- (c4);
  \draw[dep] (c2) -- (c6);
  \draw[dep] (c4) -- (c7);
  \draw[dep] (c6) -- (c7);
  \draw[dep] (c7) -- (c8);
  \draw[dep] (c8) -- (c9);
\end{tikzpicture}
\caption{\revm{Dependence of the chapters. An arrow means that the later chapter uses the earlier
one; a chapter also uses every chapter from which a chain of arrows leads to it, and arrows
implied by such chains are omitted.}}
\label{fig:roadmap}
\end{figure}
\else
\begin{figure}[ht]
\centering
\begin{tikzpicture}[
  font=\small,
  done/.style={draw, rounded corners=2pt, fill=Blue!6, text width=3.5cm, align=center,
    minimum height=1.05cm, inner sep=3pt},
  plan/.style={done, dashed, fill=none},
  dep/.style={-{Stealth[length=2mm]}, thick},
  node distance=0.9cm and 0.7cm]
  \node[done] (c1) {\ref*{sec:nets}\quad Causal indices and nets of local algebras};
  \node[done, right=of c1] (c2) {\ref*{sec:weyl}\quad The Weyl algebra};
  \node[done, right=of c2] (c3) {\ref*{sec:lattice}\quad Propagators on the lattice};
  \node[done, below=of c1] (c4) {\revb{\ref*{sec:phasespace}\quad The net of a classical phase space}};
  \node[done, below=of c2] (c6) {\revf{\ref*{sec:simple}\quad Simplicity and primitivity of the Weyl algebra}};
  \node[done, below=of c3] (c5) {\reve{\ref*{sec:latticenet}\quad The net of the lattice field}};
  \node[done, below=of c4] (c7) {\revh{\ref*{sec:green}\quad The net of a Green-hyperbolic operator}};
  \node[done, below=of c6] (c8) {\revjt{\ref*{sec:minkowski}\quad The free massless\\ field on Minkowski space}};
  \node[done, below=of c5] (c9) {\revlt{\ref*{sec:mass}\quad The Klein--Gordon\\ field on Minkowski space}};
  \draw[dep] (c1) -- (c2);
  \draw[dep] (c1.north) to[bend left=18] (c3.north);
  \draw[dep] (c2) -- (c4);
  \draw[dep, dashed] (c5) -- (c6);
  \draw[dep] (c2) -- (c6);
  \draw[dep] ([xshift=-6mm]c4.south) -- ([xshift=-6mm]c7.north);
  \draw[dep] (c3) -- (c5);
  \draw[dep] ([xshift=6mm]c4.south) -- ++(0,-0.62) -| (c5.south);
  \draw[dep] (c6) -- (c7);
  \draw[dep] (c7) -- (c8);
  \draw[dep] (c8) -- (c9);
\end{tikzpicture}
\caption{Dependence of the chapters. An arrow means that the later chapter uses the earlier
one\revi{; a chapter also uses every chapter from which a chain of arrows leads to it, and arrows
implied by such chains are omitted. Chapter~\ref{sec:simple}, for instance, rests mainly on
Chapter~\ref{sec:weyl}}. \revh{Chapter~\ref{sec:green} constructs the classical phase space from Green data, which gives
the theorem for every instance of Green data, over the derived causal index.} \revc{Chapter~\ref{sec:phasespace} uses Chapter~\ref{sec:lattice} only in its
running example, Example~\ref{ex:lattice-phase}\revo{, and in Remark~\ref{rem:massless-reversal}}.} \revf{Chapter~\ref{sec:simple} uses
Chapters~\ref{sec:phasespace} and~\ref{sec:latticenet} only in} \revg{Remark~\ref{rem:irrep} and} \revf{its last section.} \revh{Chapter~\ref{sec:green} uses \revo{Chapters~\ref{sec:lattice} and~\ref{sec:latticenet} only in remarks, motivating paragraphs and its last section}.}\revjt{ Chapter~\ref{sec:minkowski} uses Chapter~\ref{sec:green}
and, through it, Chapters~\ref{sec:phasespace} and~\ref{sec:simple}; it does not use the lattice chapters.} \revlt{Chapter~\ref{sec:mass} uses
Chapters~\ref{sec:green} and~\ref{sec:minkowski} and, through them, Chapters~\ref{sec:phasespace}
and~\ref{sec:simple}. Uses only in \revo{examples, }remarks\revo{, motivating paragraphs} or in a last section are listed
above; we draw the one from Chapter~\ref{sec:latticenet} to Chapter~\ref{sec:simple} dashed, and
follow chains of arrows along solid arrows only.}}
\label{fig:roadmap}
\end{figure}
\fi

\noindent\motiv{\rev{\revm{The main result is an abstract form of Theorem~4 of~\cite{dimock}, which is proved there for
globally hyperbolic spacetimes: whenever Green data exist (retarded and advanced Green operators
with cutoffs, Definition~\ref{def:green-data}), the net generated by the Weyl operators of localized
test functions is generating and causal, has the time-slice property for the dependence derived from
the cutoffs, and every symmetry of the Green data gives a covariance pair (Axioms 1, 3, 4 and 5;
Theorem~\ref{thm:green-axioms} and Proposition~\ref{prop:green-symmetry}). We reach the globally hyperbolic
setting of~\cite{dimock} and its form of the time-slice axiom from it on paper only
(Remarks~\ref{rem:green-continuum} and~\ref{rem:green-dep-spacetime}).} \extraonly{Chapter~\ref{sec:lattice} constructs the
propagators on the lattice\revc{, and }}\dimockonly{\revm{Chapter~\ref{sec:phasespace} proves isotony, generation,
causality and time slice for any phase-space data, and covariance for their CCR net.}}\extraonly{\revc{Chapter~\ref{sec:phasespace} proves isotony, generation,
causality and time slice for any phase-space data, and covariance for their CCR net; with
Example~\ref{ex:lattice-phase} these hold on the lattice without assumptions. \reve{Chapter~\ref{sec:latticenet}
states the lattice data as numbered results, adds covariance under the group of lattice translations
(Theorem~\ref{thm:lattice-net}), and shows that the net defined through Cauchy data is the net of a
pulled-back Weyl system, isomorphic for the CCR algebra to the lattice CCR net
(Theorem~\ref{thm:cauchy-net})}}.}}

\motiv{\rev{Primitivity, \revb{Axiom~2 of~\cite{dimock}}, is a property of the Weyl algebra alone:
\revf{for nondegenerate $\sigma$, as for the phase space of a field, Chapter~\ref{sec:simple}
proves that the algebra is simple \revg{(Theorem~\ref{thm:simple})} and deduces primitivity from the existence of an irreducible
representation (\revg{Theorems~\ref{thm:primitive} and}~\ref{thm:dimock}).}} \revh{Chapter~\ref{sec:green}
derives the phase-space data from Green operators and cutoffs (Proposition~\ref{prop:green-phase}),
which gives the axioms (Theorem~\ref{thm:green-axioms})\extraonly{, with the lattice as an instance
(Theorem~\ref{thm:lattice-green-compare})}; a continuum instance
needs such Green data, sketched \revjt{in general} (Remark~\ref{rem:green-continuum}), and\revjt{, with all
Cauchy temporal functions as cutoffs,} has the time-slice axiom in the form of~\cite{dimock}
(Remark~\ref{rem:green-dep-spacetime}).}}}

\motiv{\revjt{Chapter~\ref{sec:minkowski} constructs Green data for the free massless field on $3+1$
Minkowski space, with slabs between parallel hyperplanes as cutoffs, which gives the axioms for that
field (primitivity given an irreducible representation) with the time-slice property for those slabs (Theorem~\ref{thm:mink-green},
Remark~\ref{rem:mink-slabs}) and its covariance under the orthochronous Poincar\'e group
(Theorem~\ref{thm:mink-covariance}).}}

\motiv{\revlt{\revm{Chapter~\ref{sec:minkowski} ends by stating its construction for every local
operator with Green operators (Theorem~\ref{thm:kg-greenops}).} Chapter~\ref{sec:mass} \revm{constructs} the Green operators of the
Klein--Gordon operator $P+m^2$
for every real $m^2$ as a series in the Green operators of the wave operator, which gives the same results for that
field (Theorem~\ref{thm:kg-field}), with the massless field as the case $m^2=0$
(Corollary~\ref{cor:kg-zero})\extraonly{; Huygens' principle holds exactly when $m^2=0$
(Theorem~\ref{thm:kg-iff})}.}}
\revn{The results that correspond to~\cite{dimock} depend on few chapters: the axioms for
phase-space data (Theorem~\ref{thm:dimock}) and their covariance (Theorem~\ref{thm:ccr-cov}) use
Chapters~\ref{sec:nets}, \ref{sec:weyl}, \ref{sec:phasespace} and~\ref{sec:simple} only; \extraonly{the
lattice instance (Corollary~\ref{cor:lattice-dimock}) adds Chapters~\ref{sec:lattice}
and~\ref{sec:latticenet}; }the massless field (Theorems~\ref{thm:mink-green}
and~\ref{thm:mink-covariance}) adds Chapters~\ref{sec:green} and~\ref{sec:minkowski}, but \extraonly{neither the lattice chapters
nor }\dimockonly{not }the series of Chapter~\ref{sec:mass}.}
